\documentclass{article}
\usepackage[utf8]{inputenc}

\usepackage{fontawesome}
\usepackage{hyperref}
\usepackage{multicol}
\usepackage{booktabs}
\usepackage{dirtytalk}


\usepackage{concmath}
\usepackage[OT1]{fontenc}

\usepackage[paper=a4paper, left=1.5cm, right=1.5cm, bottom=1cm, top=1cm]{geometry}

\usepackage{longtable}

\usepackage{xcolor}
\definecolor{myGray}{RGB}{105, 93, 93}
\definecolor{darkred}{rgb}{0.55, 0.0, 0.0}
\newcommand{\shade}[1]{\textcolor{myGray}{#1}}

\hypersetup{
    colorlinks = true,
    breaklinks = true
    linkcolor  = darkred,
    urlcolor   = darkred,
}

\usepackage{amssymb}
\usepackage{enumitem}
\setlist[itemize]{label=\textcolor{darkred}{\tiny $\blacksquare$}}
\renewcommand{\labelenumi}{\textcolor{darkred}{\arabic{enumi}.}}
\renewcommand{\labelenumii}{\textcolor{darkred}{\arabic{enumi}.}}
\renewcommand{\labelenumiii}{\textcolor{darkred}{\arabic{enumi}.}}

\renewcommand\thefootnote{\textcolor{darkred}{\arabic{footnote}}}

\begin{document}
\pagestyle{empty}

\par{\centering
		{\huge Lucas Di Salvo}
	\bigskip\par}

\section*{\faUser ~~ Sobre mi}
\hrule

\
\newline
\

I'm a student and teaching assistant in the Computer Science department at the Faculty of Exact and Natural Sciences of the University of Buenos Aires (FCEN - UBA). I consider myself as a very proactive person that is interested in many topics in computer science, while also doing a deep dive on various concepts regarding programming languages. This has led me to pursue an academic career, and share this interest of mine through teaching, a job that I like very much.

\section*{\faAt ~~ Links personales}
\hrule

\
\newline
\

\begin{tabular}{l | l | l}
      \faLinkedin ~ \href{https://www.linkedin.com/in/lucas-di-salvo-6578b915a}{LinkedIn} \
     & \faGithub ~ \href{https://github.com/lucasDS-0}{GitHub}  \
     & \faEnvelope ~ \href{mailto:ldisalvo@dc.uba.ar}{ldisalvo@dc.uba.ar}
\end{tabular}

\section*{\faBook ~~ Educación}
\hrule

\
\newline
\


\footnotetext[1]{La suma de horas solo describe el tiempo dentro del aula.}
\begin{longtable}{l l}
    2017 - en curso  & \textbf{Licenciatura en Ciencias de la Computación}\\
                  & Facultad de Ciencias Exactas y Naturales, Universidad de Buenos Aires \\
					\\
                    & \textit{Tesis} \shade{(En progreso desde Enero 2026)}: Extending a reversible operational semantics for \\
                    & imperative programming languages with concurrency \\
                    & \indent $\circ$ con \href{https://dblp.org/pid/24/2104.html}{Hernán Melgratti} como director. \\
					\\
                    & \textit{Materias optativas cursadas\footnotemark[1]} (15/12 puntos
                    s, 12 necesarios y 3 puntos extra): \\
                    \\
                    & $\circ$ Tipos comportamentales y contratos \shade{(64 hs)} \\
                    & \indent \shade{Tipos sesión binarios, contratos de alto orden, tipos sesión multiparty y} \\
                    & \indent \shade{tipos sesión dependientes.} \\
                    & $\circ$ Abordaje funcional a los EDSLs \shade{(15 hs)} \hfill \shade{(\href{https://eci.dc.uba.ar/ecis-anteriores/eci2024/}{ECI2024})} \\
                    & $\circ$ The Lambda-Calculus: \\
                    & $~~$ From Simple Types to Non-Idempotent Intersection Types \shade{(15 hs)} \hfill \shade{(\href{https://eci.dc.uba.ar/ecis-anteriores/eci2024/}{ECI2024})} \\
                    & $\circ$ Reescritura, cálculo lambda y tipos \shade{(64 hs)} \\
                    & \indent \shade{Sistemas de reescritura abstracta, sistemas de reescritura de términos,} \\
                    & \indent \shade{cálculo lambda no tipado, estrategias de normalización,} \\
                    & \indent \shade{cálculo lambda simplementetipado, sustituciones explícitas.} \\
                    & $\circ$ Programación concurrente \shade{(64 hs)} \\
                    & \indent \shade{Conceptos básicos sobre concurrencia, mecanismos de sincronización,} \\
                    & \indent \shade{modelos de memoria  compartida y sus propiedades, estructuras de datos } \\
                    & \indent \shade{concurrentes,Software Transactional Memory.} \\
                    & $\circ$ Linear Logic: Syntax, Models and Extensions \shade{(15 hs)} \hfill \shade{(\href{https://eci.dc.uba.ar/ecis-anteriores/eci2023/}{ECI2023})} \\
                    & $\circ$ Modelos de programación heterogénea: \\
                    & $~~$ Portabilidad, performance y programabilidad \shade{(15 hs)} \hfill \shade{(\href{https://eci.dc.uba.ar/ecis-anteriores/eci2022/}{ECI2022})} \\
                    & $\circ$ Introducción a las APIs \shade{(24 hs)} \\
                    \\
                    & \textit{Materias obligatorias cursadas\footnotemark[1]} (20/23): \\
                    \\
                    & $\circ$ Ingeniería de Softare II \shade{(160 hs)} \\
                    & \indent \shade{Análisis dinámico de programas (algoritmos genéticos y evolutivos, fuzzing, etc.),} \\
                    & \indent \shade{análisis estático de programas (dataflow, interprocedural, intraprocedural, etc.),} \\
                    & \indent \shade{verificación de software.} \\
                    & $\circ$ Métodos Numéricos \shade{(240 hs)} \\
                    & $\circ$ Sistemas Operativos \shade{(160 hs)} \\
                    & \indent \shade{Unix y Bash, procesos y API del sistema, scheduling, sincronización y} \\
                    & \indent \shade{mecanismos deconcurrencia, manejo de memoria, drivers y I/O,} \\
                    & \indent \shade{sistemas de archivos, seguridad y virtualización.} \\
                    & $\circ$ Ingeniería de Software I \shade{(160 hs)} \\
                    & \indent \shade{OOP en SmallTalk, TDD y diseño de software.} \\
                    & $\circ$ Teoría de lenguajes \shade{(120 hs)} \\
                    & \indent \shade{Teoría de autómatas, expresiones regulares, gramáticas y parsers.} \\
                    & $\circ$ Organización del computador II \shade{(160 hs)} \\
                    & \indent \shade{Arquitecture Intel, introducción a ASM , ASM y C ABI, SIMD, microarquitectura,} \\
                    & \indent \shade{cache, paginación y segmentación, técnicas de optmización de memoria.} \\
                    & $\circ$ Paradigmas de lenguajes de programación \shade{(120 hs)} \\
                    & \indent \shade{Programación funcional en Haskell, cálculo lambda, inferencia de tipos, subtipado,}  \\
                    & \indent \shade{resolución, programación lógica en Prolog, clasificación y prototipado en Javascript.} \\
                    & $\circ$ Algoritmos y Estructuras de Datos III \shade{(240 hs)} \\
                    & \indent \shade{Teoría de grafos, backtracking, programación dinámica, algoritmos golosos,} \\
                    & \indent \shade{heurísticas y metaheurísticas, camino mínimo, MST, flujo máximo,} \\
                    & \indent \shade{introducción a la teoría de la complejidad, SAT. Hecho en C/C++ para el diseño de} \\
                    & \indent \shade{ algoritmos,y Python para la preparación de datasets, experimentación y análisis.} \\
                    & $\circ$ Lógica y Computabilidad \shade{(112 hs)} \\
                    & \indent \shade{Máquinas de Turing, Halting problem, lógica proposicional y de primer orden.} \\
                    & $\circ$ Algoritmos y Estructuras de Datos II \shade{(240 hs)} \\
                    & \indent \shade{Especificación formal de tipos abstractos de datos, árboles y estructuras de datos complejas,} \\
                    & \indent \shade{búsqueda y ordenamiento, análisis asintótico. Hecho en en C/C++.} \\
                    & $\circ$ Organización del Computador I \shade{(112 hs)} \\
                    & \indent \shade{arquitectura Von Neumann, representación de la información, diseño lógico de circuitos,} \\
                    & \indent \shade{introducción a la microarquitectura, memoria, I/O, interrupciones y buses.} \\
                    & $\circ$ Algoritmos y Estructuras de Datos I \shade{(240 hs)} \\
                    & \indent \shade{Especificación formal de programas, precondición más debil, lógica de Hoare,}  \\
                    & \indent \shade{conceptos generales de programación, testing. Hecho en C/C++.}  \\
                    & $\circ$ Probabilidad y Estadística \shade{(160 hs)} \\
                    & $\circ$ Análisis II \shade{(160 hs)} \\
                    & $\circ$ Álgebra I \shade{(208 hs)} \\
                    & $\circ$ CBC \shade{(\emph{Ciclo Básico Común})} \\
                    & \indent $*$ Análisis I, Álgebra, Introducción al pensamiento científico, Introducción al conocimiento \\
                    & \indent $~~$ de sociedad y estado, química, física. \\
                    \\
    2011 - 2016 & \textbf{Técnico en Computación} (Secundaria técnica con orientación en computación)\\
                & E.T. N°3 \say{María Sanchez de Thompson}
\end{longtable}

\section*{\faCubes ~~ Experiencia académica}
\hrule

\
\newline
\

\noindent Todos los cargos fueron obtenidos por concurso público.

\begin{longtable}{l | l}
    2024 - en curso  & \textbf{Ayudante \say{Paradigmas de Programación} at FCEN - UBA} \\
                    & Preparación y dictado de clases, junto con preparación de exámenes y su evaluación. \\
                    & \indent \shade{Main topics covered: Functional programming in Haskell, equational reasoning, natural} \\
                    & \indent \shade{ deduction (propositional logic andfirst order logic), Lambda-calculus, type inference,} \\
                    & \indent \shade{ resolution, logical programming inProlog, OOP in SmallTalk.} \\
                    & \indent $\circ$ \href{https://gestion.dc.uba.ar/media/publico/dictamenes/Dictamen_de_Concurso_Regular_de_Ayudante_de_Segunda_con_Dedicaci%C3%B3n_P_LgV1jd4.pdf}{2024 results} \shade{2-year position.} \\
                    & \indent $\circ$ \href{https://gestion.dc.uba.ar/media/publico/dictamenes/Dictamen_Concurso_de_Ayudante_de_2da_Teor%C3%ADa_2023.pdf}{2023 results} \shade{Annual position.}\\
                    \\
	2025 - 2026 (1 año) & \textbf{Teacher of \LaTeX workshop at FCEN - UBA} \\
							& Creation and lecturing of a monthly 3-hour \LaTeX workshop. \\
							& The aim of the workshop is to give tools to new students  to tackle written assignments, \\
							& reports and such, as well as being able to take stylized and detailed notes. \\
							& \indent $\circ$ \shade{This position was an experimental project. After its success and usefulness,} \\
							& \indent $~~$ \shade{it is now a regular position with public selection.} \\
                           \\
    Agosto 2023 -  & \textbf{Teaching Assistant of \say{Logic and Computability} at FCEN - UBA} \\
    Febrero 2024  & Preparation and lecturing of practical classes. Exams assessment and evaluation, \\
                   & among others. \shade{Main topics covered: Turing machines, Halting problem, propositional logic,} \\
                   & \shade{first order logic.} \\
                   & \indent $\circ$ \shade{The results platform was not in use at that time.} \\
                    \\
2020 & \textbf{Computer science popularizer at FCEN - UBA} \\
                     & Preparation of courses for high school students about different Computer Science topics, \\
                     & stands in science fairs and vocational orientation talks. \\
                     & \indent $\circ$ \shade{The results platform was not in use at that time.} \\
                     \\
    2019 - 2021     & \textbf{Programming lessons (done independently)}  \\
                    & Organization and teaching to high school students. The topics include but are not \\
                    & limited to: Python programming, OOP, structured programming, C\#, Java, \\
                    & data structures and an introduction to algorithms.
\end{longtable}


\section*{\faIndustry ~~ Experiencia en la industria}
\hrule

\
\newline
\

\begin{longtable}{l | l}
    2021 - 2022 (7 meses) & \textbf{Game Developer at \say{Globant}} \\
                           & Making of custom UIs and interactions for different game systems, \\
                           & such as a BattlePass, PopUps, Item selectors and more. \\
                           \\
    2019 (6 meses) & \textbf{Software Developer at \say{Seincomp Informática}} \\
                    &  Analysis, development and deployment of .NET (VB, C\#, Microsoft SQL Server) \\
                    & customized applications, with its maintenance and support. Web applications \\
                    & development for internal usage, with JavaScript (Boostrap), HTML, CSS.\\
\end{longtable}

\section*{\faUsers ~~ Popularización de la ciencia (Voluntariado)}
\hrule

\
\newline
\

\begin{longtable}{l | l}

Septiembre 2024 & \textbf{Science Fair - Booth Speaker} \\
Septiembre 2019 & Hosted a booth about algorithmic concepts like binary search, at the annual Computer \\
			    & science Week in my university. Made interesting through a pick and guess game, \\
			    & with high school students being the target audience. \\
			    \\
Noviembre 2020 & \textbf{Science Fair - Workshop Instructor} \\
		      & Hosted a workshop on Arduino and robotics for high school students, at the annual \\
		      & Computer Science Week in my university. This was made during Covid, so it was offered \\
		      & through Zoom and Tinkercad.\\
			  \\
Junio 2018 & \textbf{Science Fair - Guide} \\
	      & Guided high school classes through their schedule (registration, talks and venues) at the \\
	      & annual Computer Science Week in my university. \\
\end{longtable}

\section*{\faCogs ~~ Actividades de investigación}
\hrule

\
\newline
\

\begin{longtable}{l | l}
Agosto 2022 - & Asistancia a seminario de \href{https://lorel-team.github.io}{LoReL} (\textit{Lógica y reescritura para lenguajes de progrmación}) en \\
Abril 2025  & calidad de estudiante. \\
& \indent $\circ$ \shade{Los tópicos cuviertos son varios, constituidos principalmente por cálculo lambda, teoría de} \\ & \indent \shade{categorías, lógical lineal y cómputo cuántico.}
\end{longtable}


\section*{\faLanguage ~~ Lenguajes}
\hrule

\
\newline
\

\begin{longtable}{l | l}
    Español & Nativo \\
    Inglés & Avanazdo \shade{(multiples años de educación formal, aplicado tanto a la educación como a lo profesional)}\\
    Japonés & Comprensión básica \shade{(autodidacta)}
\end{longtable}

\section*{\faFileCodeO ~~ Proyectos}
\hrule
\

\begin{itemize}
    \item A project series for \textit{Algorithms and Data Structures 2} (mandatory university course), designed to develop implementations for some of the most common data structures and its functionalities (programmed in C++).
    \begin{itemize}
        \item A \href{https://github.com/lucasDS-0/Doubly_Linked_List}{Doubly Linked List}, a \href{https://github.com/lucasDS-0/BST_on_a_set}{Binary Search Tree} implemented on a \emph{set}, a \href{https://github.com/lucasDS-0/Map_on_a_trie}{Map} implemented on a \emph{trie}, and a \href{https://github.com/lucasDS-0/Priority_queue_on_a_heap}{Priority Queue} implemented on a \emph{heap}.
    \end{itemize}
    \item A small \href{https://github.com/lucasDS-0/hFetch}{tool} to fetch system information written in Haskell using System.IO to learn about monads (both in do-notation and \emph{vanilla} monadic notation).
    \item A project series for \textit{Algorithms and Data Structures 3} (mandatory university course), designed to research, analyze and develop algorithms using different programming techniques to address complex problems (programmed the core algorithms in C++, and Python3 using Jupyter notebook for the data preparation and analysis).
    \begin{itemize}
        \item An analysis on the \href{https://github.com/lucasDS-0/Subset_sum}{Subset Sum} problem. The aim of this project was to ensure the correct procedure to develop solutions by using brute force, backtracking and dynamic programming for the problem at hand, and analyze the effectiveness and efficacy of said techniques, in great detail.
        \item A project about developing heuristics over the \href{https://github.com/lucasDS-0/MIC_Polytope}{Maximum Impact Coloring Polytope} problem.
    \end{itemize}
    \item Un pequeño \href{https://github.com/lucasDS-0/ECI37_N1/tree/main}{EDSL} para la lógica proposicional con comparación aritmética. Hecho como proyeto final para la materia optativa \say{Abordaje Funcional a EDSLs}.
    \item Mi \href{https://github.com/lucasDS-0/Resume}{CV} (este documento), hecho en \LaTeX.
\end{itemize}


\end{document}
